% Chapter 2 (§2.1) signal-flow / cascade block diagram.
% Package-independent: needs only tikz + libraries
%   positioning, arrows.meta, fit, backgrounds  (all loaded in preamble.tex).
% Usage in a chapter (it is wider than \textwidth, so scale it):
%   \begin{figure}\centering\resizebox{\textwidth}{!}{% Chapter 2 (§2.1) signal-flow / cascade block diagram.
% Package-independent: needs only tikz + libraries
%   positioning, arrows.meta, fit, backgrounds  (all loaded in preamble.tex).
% Usage in a chapter (it is wider than \textwidth, so scale it):
%   \begin{figure}\centering\resizebox{\textwidth}{!}{% Chapter 2 (§2.1) signal-flow / cascade block diagram.
% Package-independent: needs only tikz + libraries
%   positioning, arrows.meta, fit, backgrounds  (all loaded in preamble.tex).
% Usage in a chapter (it is wider than \textwidth, so scale it):
%   \begin{figure}\centering\resizebox{\textwidth}{!}{% Chapter 2 (§2.1) signal-flow / cascade block diagram.
% Package-independent: needs only tikz + libraries
%   positioning, arrows.meta, fit, backgrounds  (all loaded in preamble.tex).
% Usage in a chapter (it is wider than \textwidth, so scale it):
%   \begin{figure}\centering\resizebox{\textwidth}{!}{\input{figures/fig_blockdiagram}}\caption{...}\label{fig:blockdiagram}\end{figure}
\begin{tikzpicture}[
  font=\small,
  >={Latex[length=2mm]},
  block/.style={draw, rounded corners=2pt, minimum height=1.15cm, minimum width=2.35cm,
                align=center, inner sep=3pt, fill=white},
  io/.style={draw, rounded corners=10pt, minimum height=1.15cm, minimum width=1.55cm,
             align=center, inner sep=3pt, fill=black!4},
  al/.style={font=\scriptsize, align=center},
  node distance=12mm and 16mm
]
  \node[io]                          (host)  {Host\\PC};
  \node[block, right=16mm of host]   (mcu)   {Arduino Uno\\R4 Minima};
  \node[block, right=28mm of mcu]    (drv)   {L298N H-bridge\\{\scriptsize\itshape static driver block}};
  \node[block, right=18mm of drv]    (mot)   {DC gearmotor\\{\scriptsize\itshape motor dynamic block}};
  \node[io, right=16mm of mot]       (shaft) {Output\\shaft};
  \node[block, below=18mm of shaft]  (enc)   {600\,PPR encoder\\{\scriptsize 2400 cnt/rev}};

  % forward signal path
  \draw[->] (host) -- node[al, above]{USB} (mcu);
  \draw[->] (mcu)  -- node[al, above]{PWM (D9)\\dir (D4/D5)} (drv);
  \draw[->] (drv)  -- node[al, above]{$V_{\mathrm{motor}}$} (mot);
  \draw[->] (mot)  -- (shaft);

  % feedback path
  \draw[->] (shaft) -- (enc);
  \draw[->] (enc)  -| node[al, pos=0.25, above]{A/B (D2/D3)} (mcu);

  % telemetry return
  \draw[->] (mcu)  to[bend left=22] node[al, below]{telemetry (10\,ms)} (host);

  % cascade-model highlight
  \begin{scope}[on background layer]
    \node[draw, dashed, rounded corners, fit=(drv)(mot), inner sep=11pt,
          label={[al]above:identified cascade model}] {};
  \end{scope}
\end{tikzpicture}
}\caption{...}\label{fig:blockdiagram}\end{figure}
\begin{tikzpicture}[
  font=\small,
  >={Latex[length=2mm]},
  block/.style={draw, rounded corners=2pt, minimum height=1.15cm, minimum width=2.35cm,
                align=center, inner sep=3pt, fill=white},
  io/.style={draw, rounded corners=10pt, minimum height=1.15cm, minimum width=1.55cm,
             align=center, inner sep=3pt, fill=black!4},
  al/.style={font=\scriptsize, align=center},
  node distance=12mm and 16mm
]
  \node[io]                          (host)  {Host\\PC};
  \node[block, right=16mm of host]   (mcu)   {Arduino Uno\\R4 Minima};
  \node[block, right=28mm of mcu]    (drv)   {L298N H-bridge\\{\scriptsize\itshape static driver block}};
  \node[block, right=18mm of drv]    (mot)   {DC gearmotor\\{\scriptsize\itshape motor dynamic block}};
  \node[io, right=16mm of mot]       (shaft) {Output\\shaft};
  \node[block, below=18mm of shaft]  (enc)   {600\,PPR encoder\\{\scriptsize 2400 cnt/rev}};

  % forward signal path
  \draw[->] (host) -- node[al, above]{USB} (mcu);
  \draw[->] (mcu)  -- node[al, above]{PWM (D9)\\dir (D4/D5)} (drv);
  \draw[->] (drv)  -- node[al, above]{$V_{\mathrm{motor}}$} (mot);
  \draw[->] (mot)  -- (shaft);

  % feedback path
  \draw[->] (shaft) -- (enc);
  \draw[->] (enc)  -| node[al, pos=0.25, above]{A/B (D2/D3)} (mcu);

  % telemetry return
  \draw[->] (mcu)  to[bend left=22] node[al, below]{telemetry (10\,ms)} (host);

  % cascade-model highlight
  \begin{scope}[on background layer]
    \node[draw, dashed, rounded corners, fit=(drv)(mot), inner sep=11pt,
          label={[al]above:identified cascade model}] {};
  \end{scope}
\end{tikzpicture}
}\caption{...}\label{fig:blockdiagram}\end{figure}
\begin{tikzpicture}[
  font=\small,
  >={Latex[length=2mm]},
  block/.style={draw, rounded corners=2pt, minimum height=1.15cm, minimum width=2.35cm,
                align=center, inner sep=3pt, fill=white},
  io/.style={draw, rounded corners=10pt, minimum height=1.15cm, minimum width=1.55cm,
             align=center, inner sep=3pt, fill=black!4},
  al/.style={font=\scriptsize, align=center},
  node distance=12mm and 16mm
]
  \node[io]                          (host)  {Host\\PC};
  \node[block, right=16mm of host]   (mcu)   {Arduino Uno\\R4 Minima};
  \node[block, right=28mm of mcu]    (drv)   {L298N H-bridge\\{\scriptsize\itshape static driver block}};
  \node[block, right=18mm of drv]    (mot)   {DC gearmotor\\{\scriptsize\itshape motor dynamic block}};
  \node[io, right=16mm of mot]       (shaft) {Output\\shaft};
  \node[block, below=18mm of shaft]  (enc)   {600\,PPR encoder\\{\scriptsize 2400 cnt/rev}};

  % forward signal path
  \draw[->] (host) -- node[al, above]{USB} (mcu);
  \draw[->] (mcu)  -- node[al, above]{PWM (D9)\\dir (D4/D5)} (drv);
  \draw[->] (drv)  -- node[al, above]{$V_{\mathrm{motor}}$} (mot);
  \draw[->] (mot)  -- (shaft);

  % feedback path
  \draw[->] (shaft) -- (enc);
  \draw[->] (enc)  -| node[al, pos=0.25, above]{A/B (D2/D3)} (mcu);

  % telemetry return
  \draw[->] (mcu)  to[bend left=22] node[al, below]{telemetry (10\,ms)} (host);

  % cascade-model highlight
  \begin{scope}[on background layer]
    \node[draw, dashed, rounded corners, fit=(drv)(mot), inner sep=11pt,
          label={[al]above:identified cascade model}] {};
  \end{scope}
\end{tikzpicture}
}\caption{...}\label{fig:blockdiagram}\end{figure}
\begin{tikzpicture}[
  font=\small,
  >={Latex[length=2mm]},
  block/.style={draw, rounded corners=2pt, minimum height=1.15cm, minimum width=2.35cm,
                align=center, inner sep=3pt, fill=white},
  io/.style={draw, rounded corners=10pt, minimum height=1.15cm, minimum width=1.55cm,
             align=center, inner sep=3pt, fill=black!4},
  al/.style={font=\scriptsize, align=center},
  node distance=12mm and 16mm
]
  \node[io]                          (host)  {Host\\PC};
  \node[block, right=16mm of host]   (mcu)   {Arduino Uno\\R4 Minima};
  \node[block, right=28mm of mcu]    (drv)   {L298N H-bridge\\{\scriptsize\itshape static driver block}};
  \node[block, right=18mm of drv]    (mot)   {DC gearmotor\\{\scriptsize\itshape motor dynamic block}};
  \node[io, right=16mm of mot]       (shaft) {Output\\shaft};
  \node[block, below=18mm of shaft]  (enc)   {600\,PPR encoder\\{\scriptsize 2400 cnt/rev}};

  % forward signal path
  \draw[->] (host) -- node[al, above]{USB} (mcu);
  \draw[->] (mcu)  -- node[al, above]{PWM (D9)\\dir (D4/D5)} (drv);
  \draw[->] (drv)  -- node[al, above]{$V_{\mathrm{motor}}$} (mot);
  \draw[->] (mot)  -- (shaft);

  % feedback path
  \draw[->] (shaft) -- (enc);
  \draw[->] (enc)  -| node[al, pos=0.25, above]{A/B (D2/D3)} (mcu);

  % telemetry return
  \draw[->] (mcu)  to[bend left=22] node[al, below]{telemetry (10\,ms)} (host);

  % cascade-model highlight
  \begin{scope}[on background layer]
    \node[draw, dashed, rounded corners, fit=(drv)(mot), inner sep=11pt,
          label={[al]above:identified cascade model}] {};
  \end{scope}
\end{tikzpicture}
